\documentclass[a4paper,12pt]{article}

\usepackage[utf8]{inputenc}
\usepackage[portuguese]{babel}
\usepackage{fancyhdr}
\usepackage{amsmath}
\usepackage{amsthm}
\usepackage{amssymb}
\usepackage[portuguese,algoruled,lined,linesnumbered]{algorithm2e}

% ---- PREENCHA ----
\newcommand{\Lista}{2}
\newcommand{\Nome}{Leonardo Heidi Almeida Murakami}
\newcommand{\NUSP}{11260186}
% -----------------------------

\pagestyle{fancy}
\fancyhf{}
\fancyhead[L]{\Nome}
\fancyhead[R]{NUSP \NUSP}
\fancyfoot[L]{Lista \Lista}
\fancyfoot[C]{MAC0338}

\begin{document}

\section*{Exercício 3}
\subsection*{Algoritmo}
\begin{algorithm}[H]
\SetAlgoLined
\caption{INTER}

\KwData{$k$ listas ordenadas $A_k[1..2^k], A_{k-1}[1..2^{k-1}]$ com total de $n = 2^{k+1}-1$ elementos}
\KwResult{Lista ordenada $B[1..n]$ contendo todos os $n$ elementos}

$B \leftarrow$ vetor vazio\;
\For{$i \leftarrow 0$ \KwTo k}{
    \For{$j \leftarrow 1$ \KwTo {$2^i$}}{
        $B \leftarrow A_{i}[j]$
    }
}

$q \leftarrow 1$\;
\For{$i \leftarrow 0$ \KwTo k-1}{
    $p \leftarrow 1$\;
    $r = q + 2^{i+1}$\;
    $INTERCALA(B,p,q,r)$\;
    $q = r$\;
}


\Return{$resultado$}\;
\end{algorithm}
\subsection*{Estratégia do Algoritmo}
Primeiro preenchemos o vetor B com todas as listas A, fazemos isso iterando por todas elas de 0 a k e atribuindo valor a valor para B. 
Podemos entao chamar a funcao intercala dentro do proprio vetor B, para isso, iteramos k vezes a chamada da funcao com os parametros $q$ que salvara o tamanho do bloco ja ordenado pela ultima iteracao, $p$ que sempre sera 1 para intercalarmos todo o bloco ja ordenado com o bloco novo e $r$ que é calculado através da soma de q com o tamanho da lista seguinte, ou seja, $q + 2^{i+1}$
\subsection*{Argumentar Corretude}
Sabemos que todas as listas estao, entre si, ordenadas, e que o proprio algoritmo $INTERCALA$ deixa a lista, dado que ordenada de $p$ a $q$ e de $q+1$ a $r$ ordenada. Portanto podemos ter certeza que, caso os valores de $p, q$ e $r$ forem calculados de forma correta $INTERCALA$ sempre resultara em uma lista ordenada e no final teremos a lista inteira ordenada.
\subsection*{Análise de Complexidade}
\begin{enumerate}    
    \item \textbf{Laço inicializacao:} Executamos $k$ iterações, uma para cada elemento. Em cada iteração:
    
    \begin{itemize}
        \item O laco interno roda $2^i$ vezes uma operacao $O(1)$
    \end{itemize}
    
    O custo total do laço é:
    
    \begin{equation}
    T_{\text{loop}} = \sum_{1}^{k} 2^k = 2^{k+1} - 1
    \end{equation}
    
    \item \textbf{Laço intercala:} Executamos $k$ vezes $INTERCALA$
    \begin{itemize}
        \item Na iteracao $i=1$: Unimos um bloco de tamanho 1 ($A_0$) com um de tamanho 2 ($A_1$), o custo é $O(1 + 2) = O(3)$
        \item Na iteracao $i=2$: Unimos um bloco de tamanho 3 ($A_0 + A_1$) com um de tamanho 4 ($A_2$), o custo é $O(3 + 4) = O(7)$
        \item Na iteracao $i$: Unimos um bloco de tamanho $2^i - 1$ com um de tamanho $2^i$, o custo é $O(2^{i+1} - 1)$ que, quando i=k (ou seja, na ultima iteracao), é exatamente $O(n)$
    \end{itemize}

    \item \textbf{Complexidade Total:} a complexidade total portanto é de $O(n) + O(n)$ que é $O(n)$
\end{enumerate}

\section*{Exercício 4}
\subsection*{Algoritmo}
\begin{algorithm}[H]
\SetKwProg{Def}{def}{:}{}
\SetKwRepeat{While}{while}{do}
\SetAlgoLined
\caption{INTER}

\KwData{uma lista nao ordenada $X[1..n]$}
\KwResult{uma lista ordenada $X[1..n]$}
\Def{merge-threeway(X, start, pivo1, pivo2, end)}{
    n1 = pivo1 - start + 1
    n2 = pivo2 - pivo1
    n3 = end - pivo2

    E[1..n1] $\leftarrow$ vetor vazio
    M[1..n2] $\leftarrow$ vetor vazio
    D[1..n3] $\leftarrow$ vetor vazio

    \For{i = 1 \KwTo n1}{
        E[i] = X[start + i - 1]
    }

    \For{i = 1 \KwTo n2}{
        M[i] = X[pivo1 + i]
    }

    \For{i = 1 \KwTo n3}{
        D[i] = X[pivo2 + i]
    }

    \While{}{
        
    }
}
\Def{mergesort-threeway(X, start, end)}{
    \If{$end - start$ > 2}{
        pivo1 = floor(($end+start$)/3)\;
        pivo2 = floor(2*($end+start$)/3)\;
        mergesort-threeway(X, start, pivo1)\;
        mergesort-threeway(X, pivo1+1, pivo2)\;
        mergesort-threeway(X, pivo2+1, end)\;
        merge-threeway(X, start, pivo1, pivo2, end)\;
    }
    \If{$start-end$ == 2}{
        \If{X[start] > X[end]}{
            aux = X[start]
            X[start] = X[end]
            X[end] = aux
        }
    }
}
\end{algorithm}

\end{document}